\section{Identifiability analysis}
\label{sec:identifiability}

The Maxwellian baseline calibration of Section~\ref{sec:baseline}
returns a depletion timescale $\beta = 14.27\,$ms that cannot be
reconciled with any reasonable estimate of the gas residence time.
Before drawing physical conclusions from this discrepancy, we
establish what the calibration procedure can and cannot determine
from two etch-rate anchors. We show that two-anchor calibration
constrains only one scalar combination of the model parameters,
that the choice of rate-coefficient source enters this combination
through both $\kdiss$ and the self-consistent electron density
$\nel$, and that the model's entire $r(P)$ curve at the calibration
$E/N$ is determined by two scalars regardless of how the underlying
parameters are distributed.

\subsection{Reduction to a single identifiable scalar}
\label{sec:identifiability:reduction}

Substituting \eqref{eq:power-balance}--\eqref{eq:source} into
\eqref{eq:etch} gives the explicit pipeline relation in the
isothermal limit $\alpha_T \to 0$:
\begin{equation}
  \overline{r}_{\mathrm{etch}}(P)
  \;=\;
  A_\varepsilon\,
  \frac{\kdiss\,\nel(P)\,\nSFsix}
       {1 + \beta\,\kdiss\,\nel(P)},
  \label{eq:etch-explicit}
\end{equation}
where
$A_\varepsilon \equiv \varepsilon\,\gamma_{\mathrm{wafer}}\,
(\vth/4)\,\mathcal{G}\,\nu$ collects everything proportional to
$\varepsilon$. It is convenient to absorb the rate-source
dependence into a single power-independent quantity
\begin{equation}
  \alpha
  \;\equiv\;
  \frac{\kdiss\,\nel(P)}{P}
  \;=\;
  \frac{\kdiss}
       {\nSFsix\,V\,\varepsilon_{\mathrm{loss}}\,q_e},
  \label{eq:alpha}
\end{equation}
which is independent of $P$ at fixed $E/N$ because the power balance
\eqref{eq:power-balance} makes $\nel$ linear in $P$. The depletion
factor reduces to $1 + \beta\alpha P$, and
\eqref{eq:etch-explicit} becomes
\begin{equation}
  \overline{r}_{\mathrm{etch}}(P)
  \;=\;
  A_\varepsilon\,\nSFsix\,\alpha\,
  \frac{P}{1 + \beta\alpha\,P}.
  \label{eq:etch-alpha}
\end{equation}
Equation~\eqref{eq:etch-alpha} expresses the entire $r(P)$ curve as
a function of exactly two effective scalars: the prefactor
$A_\varepsilon\,\nSFsix\,\alpha$ and the depletion product
$\beta\alpha$.

Taking the ratio of \eqref{eq:etch-alpha} at the two anchors
cancels the prefactor:
\begin{equation}
  \frac{r_H}{r_L}
  \;=\;
  \frac{P_H}{P_L}\cdot
  \frac{1 + \beta\alpha\,P_L}{1 + \beta\alpha\,P_H}.
  \label{eq:ratio-constraint}
\end{equation}
This is one scalar equation in the single unknown $\beta\alpha$.
Cross-multiplying and collecting terms yields the closed-form
solution
\begin{equation}
  \beta\alpha
  \;=\;
  \frac{P_H\,r_L - P_L\,r_H}{P_H\,P_L\,(r_H - r_L)}.
  \label{eq:beta-alpha}
\end{equation}
Substituting the anchor values
$(P_L,r_L) = (200,1400)$ and $(P_H,r_H) = (2000,2270)$ in units of
$(\mathrm{W},\mathrm{nm\,min^{-1}})$ gives
$\beta\alpha = 6.7414\times10^{-3}\,$W$^{-1}$. Once $\beta\alpha$
is fixed, the prefactor in \eqref{eq:etch-alpha} is determined by
either anchor in one division,
\begin{equation}
  A_\varepsilon\,\nSFsix\,\alpha
  \;=\;
  \frac{r_L\,(1 + \beta\alpha\,P_L)}{P_L},
  \label{eq:A-eps}
\end{equation}
and equations~\eqref{eq:beta-alpha}--\eqref{eq:A-eps} together
exhaust the information content of the two-anchor calibration.

The implication is that the model's entire predicted curve $r(P)$
at the calibration $E/N$ is parameterized by exactly two scalars:
$\beta\alpha$ and $A_\varepsilon\nSFsix\alpha$. The underlying
parameters $(\varepsilon, \beta, \kdiss,
\varepsilon_{\mathrm{loss}})$ are not individually identifiable;
they live on a continuous family of $(\varepsilon,\beta)$ pairs
that all give the same $r(P)$ curve as long as $\alpha$ is rescaled
accordingly. The fitted $\beta$ is therefore not directly a
physical quantity; it is the value required for the depletion
product $\beta\alpha$ to match the anchor data, given a specific
assumed value of $\alpha$. Different assumptions about the rate
source yield different fitted $\beta$, even though the predicted
$r(P)$ curve at the calibration $E/N$ is identical in every case.

\subsection{Heated case and the numerical invariant}
\label{sec:identifiability:heated}

The derivation of
Section~\ref{sec:identifiability:reduction} is exact in the
isothermal limit $\alpha_T \to 0$. Under the linear gas-heating
model $T_g(P) = T_0 + \alpha_T P$ used in the baseline calibration,
several additional factors in equation~\eqref{eq:etch} acquire
power dependence: the F-atom thermal speed
$\vth \propto T_g^{1/2}$, the diffusion coefficient
$D \propto T_g^{3/2}$ through $\mathcal{G}(D)$, and the ideal-gas
neutral density $\nSFsix \propto T_g^{-1}$. The simple Hill form
no longer describes $r(P)$ exactly, and
\eqref{eq:beta-alpha} no longer gives the fitted $\beta\alpha$ in
closed form.

The identifiability \emph{structure} is nevertheless preserved: in
the heated case there is still a single scalar combination of the
underlying parameters that the two anchors determine, along with a
magnitude prefactor that $\varepsilon$ absorbs. We extract the
heated-case invariant numerically as the dimensionless product
\begin{equation}
  \mathcal{I}
  \;\equiv\;
  \beta\,\kdiss\,\nel(P_H),
  \label{eq:invariant-full}
\end{equation}
where $\nel(P_H)$ is computed self-consistently from the power
balance at the high-power anchor using the rate-coefficient source
under test. For the Maxwellian baseline,
$\beta = 14.27$\,ms,
$\kdiss = 1.109\times10^{-17}$\,m$^3$\,s$^{-1}$,
$\nel(P_H) = 1.574\times10^{19}$\,m$^{-3}$ at
$P_H = 2000\,$W, giving $\mathcal{I}_{\mathrm{Max}} = 2.491$. The
rate-source-dependent consequence of the invariant constraint is
expressed as
\begin{equation}
  \frac{\beta_{\mathrm{new}}}{\beta_{\mathrm{old}}}
  \;=\;
  \frac{\kdiss^{\mathrm{old}}}{\kdiss^{\mathrm{new}}}
  \;\cdot\;
  \frac{\nel^{\mathrm{old}}(P_H)}{\nel^{\mathrm{new}}(P_H)},
  \label{eq:beta-shift}
\end{equation}
which holds exactly because $\mathcal{I}$ is preserved by
construction. Section~\ref{sec:boltzmann} uses
\eqref{eq:beta-shift} to make a falsifiable prediction about the
shift in fitted $\beta$ when the Maxwellian rate source is replaced
with a two-term Boltzmann rate source, and verifies the prediction
numerically.

\subsection{Numerical verification}
\label{sec:identifiability:numerics}

To verify \eqref{eq:beta-alpha}, we disable gas heating
($\alpha_T = 0$), keep the Maxwellian rate-coefficient backend
fixed, and fit $(\varepsilon,\beta)$ against the anchors. The
closed form predicts
$\beta\alpha = 6.7414\times10^{-3}\,$W$^{-1}$, which, using the
power-balance value of $\alpha$ at the isothermal $E/N$, gives
$\beta = 77.27\,$ms. A direct numerical fit in the isothermal
limit returns $\beta = 77.26\,$ms, in agreement with the closed
form to better than $0.02\,\%$.

As a further check, we rescale $\kdiss$ artificially by a factor
$f$ (still isothermal) and refit $(\varepsilon,\beta)$. The
identifiability theorem predicts $\beta \propto 1/f$ while every
observable prediction remains unchanged.
Table~\ref{tab:scaling} collects the result for four values of $f$.

\begin{table}[t]
\centering
\caption{Numerical verification of the isothermal identifiability
theorem under artificial $\kdiss$ rescaling. The fitted depletion
timescale $\beta$ scales as $1/f$ to within $0.02\,\%$, and the
predicted etch-rate curves are indistinguishable to the tolerance
of the numerical fit across the four rescaled fits.}
\label{tab:scaling}
\begin{tabular}{cccc}
\toprule
$f$ & $\beta_{\mathrm{fit}}$ (ms) &
$f\cdot\beta_{\mathrm{fit}}$ (ms) & predicted slope \\
\midrule
$1.00$ & $77.27$  & $77.27$ & --- \\
$0.50$ & $154.53$ & $77.27$ & $1.000$ \\
$0.30$ & $257.57$ & $77.27$ & $1.000$ \\
$0.10$ & $772.70$ & $77.27$ & $1.000$ \\
\bottomrule
\end{tabular}
\end{table}

The product $f\beta_{\mathrm{fit}}$ is constant to numerical
tolerance, confirming that the identifiable scalar is
$\beta\kdiss$ when $\varepsilon_{\mathrm{loss}}$ is held fixed by
holding the rate source fixed.

\subsection{Physical interpretation}
\label{sec:identifiability:physics}

The identifiability analysis clarifies what the depletion
calibration actually measures. In the isothermal case the two
anchors constrain $\beta\alpha$, the dimensionless ratio of the
electron-impact dissociation rate to the pumping rate. In the
heated case they constrain the equivalent numerical invariant
$\mathcal{I} = \beta\,\kdiss\,\nel(P_H)$. Any model with the
correct value of this invariant and the correct prefactor will
reproduce the anchor etch rates and the entire intermediate $r(P)$
curve. The individual values of $\beta$, $\kdiss$, $\nel$, and
$\varepsilon$ are not constrained by the data on their own merits;
they are constrained jointly through their participation in the
identifiable product.

This explains the apparent failure of the Maxwellian baseline
calibration. The Maxwellian assumption overestimates $\kdiss$ in
\sfsix{} at $E/N = 100\,$Td by a factor of approximately 29 owing
to the depletion of the high-energy EEDF tail by attachment in the
real two-term Boltzmann distribution. A larger assumed $\kdiss$
forces $\beta$ to take a smaller value to keep $\mathcal{I}$ on the
locus determined by the anchors, and the resulting
$\beta = 14\,$ms is the value the model requires to reproduce the
data given an erroneous $\kdiss$. The small-$\beta$ value is not a
failure of the depletion mechanism nor a sign that the depletion
model is inappropriate; it is the arithmetic consequence of an EEDF
error being absorbed into the single fitted parameter that
participates in the identifiable product alongside $\kdiss$.

The methodological lesson is that the calibration result
$\beta = 14\,$ms cannot be interpreted as a residence time in the
first place. Two etch-rate measurements at one
$(E/N,x_{\mathrm{Ar}})$ contain only enough information to
determine one scalar combination of the depletion and rate
parameters, and $\beta$ on its own is not that combination. Either
an independent estimate of $\kdiss$, or a measurement at a
different operating condition that breaks the at-anchor degeneracy
(Section~\ref{sec:observability}), is required before the fitted
depletion parameter can be assigned a physical interpretation.
